\bookchapter{Capitolo 39: Ingegneria del contesto — far capire all'AI le tue intenzioni}

\chapterintro{Ti è mai capitato che aprendo una nuova finestra di chat l'AI sembrasse amnesica e dimenticasse tutto del tuo progetto? Non lo fa apposta: è «smemorata» per natura. Questo capitolo ti insegna a installarle un «sistema di memoria a lungo termine» e ti presenta uno dei paradigmi di sviluppo più caldi del 2026: lo sviluppo guidato dalle specifiche.}

\booksection{Il file di memoria di progetto: la «memoria a lungo termine» dell'AI}

I modelli AI hanno un limite fondamentale: la memoria è «a breve termine». A ogni nuova conversazione l'AI sa zero del tuo progetto. Devi ripeterle «usiamo React», «il backend è Node.js», «il database è PostgreSQL», «lo stile del codice usa doppi apici». Non solo inefficiente, ma facile da dimenticare — come formare ogni giorno un nuovo collega.

La soluzione è il «file di memoria di progetto»: un file di testo nella root del progetto che gli strumenti AI leggono automaticamente a ogni conversazione. Ogni strumento ha la propria convenzione di nome:

\begin{center}
\begin{tabularx}{\linewidth}{|X|X|}
\hline
\textbf{Strumento} & \textbf{Nome file} \tabularnewline
\hline
Cursor & \texttt{.\allowbreak cursorrules} \tabularnewline
Claude Code & \texttt{CLAUDE.\allowbreak md} \tabularnewline
GitHub Copilot & \texttt{copilot-\allowbreak instructions.\allowbreak md} \tabularnewline
Soluzione generica & \texttt{AGENTS.\allowbreak md} \tabularnewline
\hline
\end{tabularx}
\end{center}

\begin{center}\small\color{gray}Convenzioni di nome del file di memoria per ogni strumento\end{center}

Nomi diversi, stesso principio: scrivi in quel file ciò che l'AI deve «ricordare sempre». Un buon file di memoria di progetto contiene: panoramica del progetto (che cos'è), stack tecnologico (quali tecnologie), standard di codice (naming, indentazione, commenti), struttura delle cartelle (come è organizzato il codice), convenzioni di sviluppo (strategia dei branch, commit, test). Il CLAUDE.md di Karpathy ha solo 70 righe ma copre tutte le informazioni chiave: conciso e completo.

\begin{lstlisting}
# Progetto: app personale di gestione finanze
## Panoramica
Un'app full-stack di finanza personale che aiuta a registrare entrate e uscite quotidiane.
## Stack tecnologico
- Frontend: React + TypeScript + Tailwind CSS
- Backend: Node.js + Express + TypeScript
- Database: SQLite (sviluppo) / PostgreSQL (produzione)
- Test: Vitest
## Convenzioni del codice
- Usa virgolette doppie
- Indentazione di 2 spazi
- Nomi delle funzioni in camelCase
- Definizioni di tipo in PascalCase
- Ogni funzione deve avere un commento JSDoc
## Struttura delle cartelle
/src
  /frontend  - codice frontend
  /backend   - codice backend
  /shared    - tipi condivisi tra frontend e backend
  /tests     - file di test
## Regole di sviluppo
- Messaggi di commit in italiano
- Le nuove funzionalità richiedono prima i test
- Le API rispondono sempre con { success, data, error }
\end{lstlisting}

\begin{tcbwarn}{Attenzione: non trasformarlo in un «secondo documento di requisiti»}
Fai attenzione alla lunghezza: non deve diventare un documento di requisiti di decine di pagine. Lo scopo è «far costruire rapidamente il contesto all'AI», non «descrivere in dettaglio ogni funzione». Mantienilo conciso, strutturato, solo «ciò che l'AI deve sapere».

\end{tcbwarn}

\booksection{Sviluppo guidato dalle specifiche (SDD)}

Lo Spec-Driven Development (SDD) è un paradigma di sviluppo maturato progressivamente nel 2026: prima le specifiche, poi il codice. A differenza del TDD tradizionale (test-driven development), le «specifiche» dell'SDD non sono solo casi di test, ma anche descrizioni dei requisiti, definizioni di interfacce, modelli di dati. L'AI genera il codice sulla base di queste specifiche; agli umani spetta scriverle e mantenerle.

Il flusso dell'SDD in cinque passi:

\begin{enumerate}
  \item Descrivi i requisiti in linguaggio naturale (in un file .md);
  \item Definisci interfacce e strutture dati (nel file di specifiche);
  \item Lascia che l'AI generi il codice dalle specifiche;
  \item Esegui i test per verificare la conformità alle specifiche;
  \item Se non è conforme, correggi le specifiche o fai correggere il codice all'AI.
\end{enumerate}

In questo flusso il file di specifiche è la «fonte di verità unica»: il codice si può rigenerare, le specifiche devono essere giuste prima.

Scrivi il file di specifiche

Le specifiche sono complete?

Completa

L'IA genera codice dalle specifiche

Esegui i test per verificare l'output

Incompleta / test falliti

Problema nel codice correggi il codice; problema nelle specifiche correggi le specifiche

Torna alla verifica

Tutto superato consegna

\begin{center}\small\color{gray}Figura 3: flusso di lavoro dello sviluppo guidato dalle specifiche (SDD)\end{center}

Il vantaggio dell'SDD: le specifiche sono più stabili del codice. Quando cambiano modello o strumento, le tue specifiche restano valide e basta rigenerare il codice con il nuovo strumento AI. Ecco perché il libro insiste che «la metodologia conta più degli strumenti»: il file di specifiche è il veicolo concreto della tua metodologia.

\booksection{Gestione della finestra di contesto}

Anche con il file di memoria di progetto, in chat può emergere il problema del «contesto insufficiente». La finestra di contesto dei modelli AI (la quantità di informazioni elaborabili contemporaneamente) è limitata. Quando la conversazione è troppo lunga e i file citati troppi, l'AI «dimentica» le informazioni iniziali e la qualità cala — come una chiacchierata di sei ore con un amico che ha già dimenticato cosa gli avevi detto all'inizio.

La strategia centrale è il «riferimento a livelli»:

\begin{center}
\begin{tabularx}{\linewidth}{|X|X|X|}
\hline
\textbf{Livello di contesto} & \textbf{Contenuto} & \textbf{Gestione} \tabularnewline
\hline
Livello 1: memoria di progetto & Stack, specifiche, struttura cartelle & Scritto in file come CLAUDE.md, caricato automaticamente, senza ripetizioni \tabularnewline
Livello 2: riferimenti a file & I file di codice coinvolti nell'attività corrente & Riferisci con il simbolo @ in modo preciso, non incollare interi file in chat \tabularnewline
Livello 3: contesto di dialogo & Cronologia della conversazione sull'attività corrente & A fine attività apri una nuova chat per evitare che il contesto si gonfi \tabularnewline
Livello 4: istruzioni immediate & Requisiti specifici di questo passaggio & Chiari, concisi, mirati; evita descrizioni prolisse \tabularnewline
\hline
\end{tabularx}
\end{center}

\begin{center}\small\color{gray}Strategia di gestione a livelli della finestra di contesto\end{center}

\booksection{I punti chiave della pratica dell'ingegneria del contesto}

L'ingegneria del contesto non è una configurazione da fare una volta sola, ma un processo di ottimizzazione continua. Ecco alcuni punti chiave verificati più e più volte nella pratica:

\begin{enumerate}
  \item Tieni i file di memoria aggiornati insieme all'evoluzione del progetto. Quando introduci una nuova tecnologia o cambi l'architettura, aggiorna subito il file di memoria invece di spiegarlo al volo nella conversazione.
  \item Usa bene i riferimenti @ invece di incollare testi interi. La maggior parte degli strumenti di programmazione AI supporta la sintassi @file o @folder, che inietta in background il contenuto del file nel contesto dell'AI: è più efficiente dell'incollaggio manuale e non inquina la conversazione.
  \item Le conversazioni devono essere brevi e incisive. Una conversazione gestisce un solo compito o una sola funzionalità. Finito? Apri una nuova conversazione per il successivo. Le conversazioni lunghe non solo consumano la finestra di contesto, ma fanno anche “distrarre” l'AI.
  \item Fai “vedere” all'AI i risultati invece di “descriverli”. Quando incontri un bug, cita direttamente il log degli errori con @ invece di raccontarlo a parole tue con un vago “c'è tipo un errore”. Per l'AI è molto più efficace vedere le informazioni originali che la tua parafrasi.
  \item Pulizia e archiviazione periodiche. Le conversazioni concluse e i file temporanei occupano spazio di lavoro. Pulire di tanto in tanto ciò che non serve e tenere ordinata la directory del progetto aiuta anche l'AI a capire la struttura del progetto in modo più efficiente.
\end{enumerate}

\begin{tcbpractice}{Mettiti alla prova: scrivi un file di memoria di 70 righe per il tuo progetto}
Scegli un progetto che stai facendo (o che intendi fare) e scrivi un CLAUDE.md: panoramica del progetto (2-3 frasi), stack tecnologico (elencalo), regole del codice (3-5 punti), struttura delle directory (disegnala ad albero), convenzioni di sviluppo (2-3 punti). Resta entro 70 righe: quando avrai finito scoprirai che anche la tua comprensione del progetto è diventata molto più chiara.

\end{tcbpractice}

\begin{tcbsummary}{Riepilogo del capitolo}
\begin{itemize}
  \item La memoria dell'AI è “a breve termine”: a ogni nuova conversazione la sua conoscenza del progetto riparte da zero
  \item I file di memoria del progetto (CLAUDE.md e simili) sono la “memoria a lungo termine” dell'AI, caricati automaticamente a ogni conversazione
  \item Un buon file di memoria = panoramica + stack + regole + directory + convenzioni: bastano 70 righe
  \item SDD, sviluppo guidato dalle specifiche: prima la specifica poi il codice; la specifica è la “fonte unica di verità”, più stabile del codice
  \item Il contesto si gestisce su quattro livelli: memoria del progetto riferimenti @file cronologia della conversazione istruzioni immediate
  \item Cinque punti chiave pratici: aggiornamenti sincronizzati, riferimenti @, conversazioni brevi, fornire le informazioni originali, pulizia periodica
\end{itemize}

\end{tcbsummary}
